% ====================================================================
%  KANDA -- IEEE Conference Paper (Phase 4)
%  Compile: pdflatex -> pdflatex (x2 if bib added later)
% ====================================================================
\documentclass[conference]{IEEEtran}
\IEEEoverridecommandlockouts

\usepackage{cite}
\usepackage{amsmath,amssymb,amsfonts}
\usepackage{newtxtext,newtxmath}
\usepackage{algorithm}
\usepackage{algorithmic}
\usepackage{graphicx}
\graphicspath{{./}{../images/}}
\usepackage{textcomp}
\usepackage{booktabs}
\usepackage{array}
\usepackage{tabularx}
\usepackage{multirow}
\usepackage{enumitem}
\usepackage{xcolor}
\usepackage{listings}
\usepackage{float}
\usepackage{microtype}
\usepackage[hidelinks]{hyperref}

% Fixed-size figure slot, aspect ratio preserved (#1=width, #2=height)
\newcommand{\slotfig}[3]{%
  \begin{minipage}[c][#2][c]{#1}%
    \centering
    \includegraphics[width=\linewidth,height=#2,keepaspectratio]{#3}%
  \end{minipage}%
}

% Full-width tables in IEEE single-column floats
\newcolumntype{Y}{>{\raggedright\arraybackslash}X}
\newcolumntype{C}{>{\centering\arraybackslash}X}
\newcolumntype{R}{>{\raggedleft\arraybackslash}X}

\lstset{
  basicstyle=\ttfamily\scriptsize,
  breaklines=true,
  breakatwhitespace=false,
  breakautoindent=true,
  postbreak=\mbox{\textcolor{gray}{$\hookrightarrow$}\space},
  frame=single,
  numbers=none,
  captionpos=b,
  aboveskip=6pt,
  belowskip=2pt,
  belowcaptionskip=2pt,
  keepspaces=true,
  columns=fullflexible,
}

\def\BibTeX{{\rm B\kern-.05em{\sc i\kern-.025em b}\kern-.08em
    T\kern-.1667em\lower.7ex\hbox{E}\kern-.125emX}}

\newcommand{\safe}{\texttt{stop}}
\newcommand{\groq}{\textit{Groq Llama 3.3}}
\newcommand{\nvidia}{\textit{NVIDIA NIM Llama 3.2 Vision}}

% ====================================================================
\begin{document}

\title{KANDA: A Multimodal Embodied Robot Agent Powered by\\LLMs with Hardware-Aware Action Generation}

\author{%
\IEEEauthorblockN{Sudarshan T S}
\IEEEauthorblockA{%
\textit{Dept.\ of Information Science and Engineering}\\
\textit{RV College of Engineering (Autonomous)}\\
Bengaluru 560059, India\\
\texttt{sudarshants.sit24@rvce.edu.in}}
\and
\IEEEauthorblockN{Sushmitha N}
\IEEEauthorblockA{%
\textit{Dept.\ of Information Science and Engineering}\\
\textit{RV College of Engineering (Autonomous)}\\
Bengaluru 560059, India\\
\texttt{sushmithan@rvce.edu.in}}
}

\maketitle

% ====================================================================
\begin{abstract}
We present \textbf{KANDA}, a multimodal embodied agent built for under \$77 from an ESP32, a Raspberry Pi~4, and a camera module.
The robot accepts spoken and text commands, navigates indoors, and searches for named objects using \groq{} for text reasoning and \nvidia{} for visual understanding.
Reliability rests on three orthogonal layers aligned with Brooks-style subsumption: (i)~a body-context prompt that grounds the model in live sensors and action limits, (ii)~a deterministic \emph{critic} validator on the Pi that maps any LLM output to a legal command or \safe{}, and (iii)~a firmware \emph{reflex} on the ESP32 that stops forward motion when the front ultrasonic range drops below 15\,cm, independently of cloud latency.
Ablation over 10 trials per condition shows distinct failure modes when each layer is removed (e.g., 15\% unsafe plans without scene context; 60\% redundant search steps without episodic memory; 30\% near-miss events without firmware stop).
On 60 voice utterances, macro-averaged intent F1 is 0.94; visual search succeeds in 8/10 trials; eight injected faults all terminate in a safe stop state.
\end{abstract}

\begin{IEEEkeywords}
embodied AI, deliberative--critic--reflex architecture, vision-language grounding, command validation, episodic visual search, layered safety, ESP32, Raspberry Pi
\end{IEEEkeywords}

% ====================================================================
\section{Introduction}
\label{sec:intro}
% ====================================================================

Recent embodied AI systems rely on costly hardware and large training corpora: PaLM-E~\cite{driess2023palme} uses a 562B-parameter model; SayCan~\cite{ahn2022saycan} targets mobile manipulators priced above \$100k; RT-1~\cite{brohan2022rt1} and RT-2~\cite{brohan2023rt2} depend on large teleoperation and web-scale fine-tuning datasets.
This work asks whether comparable \emph{behaviours} (voice interaction, indoor navigation, object search, collision avoidance) can run on commodity parts (ESP32, Pi~4, Pi Camera) with zero robot-specific training and a hosted vision-language API.

Four engineering gaps motivate the design:
\textbf{(1) No body context.} Without explicit movement primitives, speed limits, and sensor readings, the text model emits infeasible commands, as reported on larger platforms~\cite{ahn2022saycan,vemprala2024chatgpt}.
\textbf{(2) Intent ambiguity.} A motion command requires a single validated action; a search command requires a multi-step vision loop inspired by observation-driven planners~\cite{yao2023react,huang2022innermonologue}.
\textbf{(3) Model garbage.} JSON plans still arrive malformed, wrapped in markdown, or with hallucinated action names~\cite{huang2023grounded}.
\textbf{(4) Timescale mismatch.} Cloud inference takes 1.5--4\,s while obstacles can appear within 100\,ms~\cite{brooks1986robust}.

KANDA addresses these with a body-context assembler, intent routing, a total validator function $V$, and firmware-level reflexes.

\begin{figure}[t]
\centering
\slotfig{0.9\columnwidth}{4.2cm}{kanda_robot}
\caption{KANDA hardware: ESP32 DevKit, Raspberry Pi~4, Pi Camera v2.1, three HC-SR04 ultrasonic sensors, SSD1306 OLED, and TB6612FNG motor driver.}
\label{fig:proto}
\end{figure}

Contributions: (1)~a \$77 multimodal agent (voice, vision, text/Telegram); (2)~body-context grounding with shared action vocabulary; (3)~intent classification plus episodic visual search; (4)~deliberative--critic--reflex safety with ablation evidence; (5)~a seven-state machine that preserves fast obstacle handling despite slow cloud calls; (6)~demonstration and remote-control modes.
Implementation artifacts are described in the project repository (vision module, firmware, and test scripts).

% ====================================================================
\section{Related Work}
\label{sec:related}
% ====================================================================

SayCan~\cite{ahn2022saycan} learns affordances from demonstrations; KANDA encodes affordances in prompts (zero training).
Huang et al.~\cite{huang2022zeroshot} and chain-of-thought prompting~\cite{wei2022cot} show that careful framing improves executability; our body-context extends this with live telemetry.
Socratic Models~\cite{zeng2023socratic}, Code as Policies~\cite{liang2023code}, and ProgPrompt~\cite{singh2023progprompt} compose language into control programs; we restrict the interface to JSON lines validated before UART transmission.
Inner Monologue~\cite{huang2022innermonologue} and ReAct~\cite{yao2023react} interleave observations with decisions; our search loop uses VLM scene descriptions and episodic deduplication rather than explicit chain-of-thought text from the robot.

For safety, Amodei et al.~\cite{amodei2016safety} argue for structural safeguards beyond model scaling.
Brooks~\cite{brooks1986robust} places fast reactive layers above slow deliberation; the Dynamic Window Approach~\cite{fox1997dwa} informs ultrasonic thresholds.
Edge and fog computing surveys~\cite{bonomi2012fog,shi2016edge} motivate executing safety-critical logic on the microcontroller while keeping reasoning on the Pi.
ProgPrompt and similar systems implicitly filter programs; KANDA makes the critic explicit and pairs it with a hardware reflex that cannot be bypassed by prompt engineering alone.

% ====================================================================
\section{Formal Problem Definition}
\label{sec:formal}
% ====================================================================

The action vocabulary is $\mathcal{A} = \{\text{forward}, \text{backward}, \text{left}, \text{right}, \text{slight\_left}, \text{slight\_right}, \text{stop}\}$ with speed in $[0, 255]$.

\textbf{Definition 1 (Safe command).}
$c = (\text{action}, \text{speed})$ is \emph{safe} iff Eq.~\eqref{eq:safety} holds:
\begin{equation}
\label{eq:safety}
c.\text{action} \in \mathcal{A} \;\wedge\; c.\text{speed} \in \mathbb{Z} \;\wedge\; 0 \leq c.\text{speed} \leq 255
\end{equation}

\textbf{Definition 2 (Critic / validator).}
$V:\text{Any} \to \mathcal{A} \times [0,255]$ must satisfy output safety and totality:
\begin{equation}
\label{eq:validator}
\forall\, x:\; V(x) \text{ satisfies Eq.~\eqref{eq:safety}} \;\wedge\; V \text{ terminates without exception}
\end{equation}
Invalid actions map to $\safe{}$; non-integer speeds map to $\safe{}$; numeric speeds outside range are \emph{clamped} to $[0,255]$ before transmission (Listing~\ref{lst:validator}).

\textbf{Definition 3 (Timing constraint).}
Measured reflex latency $t_{\text{reflex}}$ (firmware stop after forward motion into $<$15\,cm) must remain well below cloud latency $t_{\text{cloud}}$:
\begin{equation}
\label{eq:timing}
t_{\text{reflex}} \ll t_{\text{cloud}}, \quad
t_{\text{reflex}} \approx 47\,\text{ms} \;\text{(mean)}, \;
t_{\text{cloud}} \in [1500, 4000]\,\text{ms}
\end{equation}
The ESP32 runs a single cooperative loop at $\sim$10\,Hz (\texttt{delay(100)}); reflex logic executes in that loop without awaiting cloud inference.

\textbf{Definition 4 (Episodic similarity).}
For scene strings $s_i$, $s_j$, let $\text{sim}(s_i,s_j)$ be the Ratcliff--Obershelp ratio implemented by Python \texttt{difflib.SequenceMatcher.ratio}~\cite{ratcliff1988}:
\begin{equation}
\label{eq:similarity}
\text{sim}(s_i, s_j) \in [0,1], \quad
\text{visited if } \text{sim}(s_{\text{new}}, s_{\text{mem}}) > \tau
\end{equation}
with $\tau{=}0.7$, lowered to $0.5$ after five consecutive visits to the same scene class.

\textbf{Definition 5 (Grounding completeness).}
Let $\mathcal{P}$, $\mathcal{V}$, $\mathcal{F}$ be actions in the prompt, validator/executor whitelist, and firmware parser, respectively.
\begin{equation}
\label{eq:grounding}
\mathcal{P} = \mathcal{V} = \mathcal{F} = \mathcal{A}
\end{equation}
The vision module enforces alignment through shared move-action sets in \texttt{body\_context.py}, \texttt{plan\_executor.py}, and \texttt{ai\_layer/config.py}.

% ====================================================================
\section{System Design}
\label{sec:arch}
% ====================================================================

\subsection{Three Tiers}

Three tiers (Fig.~\ref{fig:arch}) each reach a safe state if upper tiers fail.

\begin{figure}[t]
\centering
\slotfig{0.9\columnwidth}{5.0cm}{architecture}
\caption{Three-tier layout: cloud deliberation, edge critic and state machine, device reflex and sensing.}
\label{fig:arch}
\end{figure}

The \textbf{cloud} tier (Groq for text, NVIDIA NIM for vision; temperature 0.1 for classification) performs intent labeling, JSON planning, and scene/target analysis; latency is 1.5--4\,s and outputs are treated as untrusted.
The \textbf{edge} tier (Raspberry Pi~4, Python vision module) hosts the state machine, body-context builder, plan executor, openWakeWord~\cite{openwakeword2024} wake phrase \emph{Hey Kanda}, gTTS speech output, and Telegram polling.
The \textbf{device} tier samples three HC-SR04 sensors each loop, parses one JSON command per line, and applies an emergency stop when moving forward with front range $<$15\,cm.

\subsection{Deliberative, Critic, and Reflex Layers}
\label{sec:dcr}

\textbf{Deliberative (cloud + prompt).} The text model proposes actions and plans conditioned on body context (Listing~\ref{lst:bodyctx}). The vision model answers "is target visible?" queries on captured frames.
\textbf{Critic (edge).} $V$ rejects unknown actions and non-numeric speeds; clamps otherwise valid speeds (Listing~\ref{lst:validator}).
\textbf{Reflex (device + edge cancel).} Firmware stops motors and reports \texttt{OBSTACLE} in telemetry; the Pi heartbeat thread sets a cancel event so THINKING/SEARCHING states abort without blocking on cloud replies.
This split matches Brooks subsumption: reflex handles physics; critic handles syntax; deliberation handles semantics.

\subsection{Interaction Data-Flow}

Fig.~\ref{fig:dataflow} traces ``go forward'' from microphone to motors.

\begin{figure*}[t]
\centering
\slotfig{0.92\textwidth}{4.1cm}{dataflow}
\caption{Interaction data-flow from microphone through wake word, speech-to-text, intent and planning, critic validation, UART, and ESP32 execution.}
\label{fig:dataflow}
\end{figure*}

\subsection{Wire Protocol}

The Pi sends one JSON object per line at 115{,}200 baud, e.g.,
{\scriptsize\texttt{\{"action":"forward","speed":120,"state":"searching"\}}}.
The ESP32 replies with newline-delimited telemetry, e.g.,
{\scriptsize\texttt{F:30.5 L:21.0 R:16.4 -> FORWARD}}.
Text framing enables debugging with standard serial monitors.

% ====================================================================
\section{Mechanisms}
\label{sec:methodology}
% ====================================================================

\subsection{Body-Context Assembler}
\label{sec:bodyctx}

Every model call is prefixed with capabilities, limits, live ranges, optional warnings, scene text, and recent actions (Listing~\ref{lst:bodyctx}).

\begin{lstlisting}[caption={Body-context block (rebuilt each inference call)},
  label={lst:bodyctx}, float=tp]
You are KANDA, a two-wheeled indoor robot.
Capabilities: forward, backward, left, right,
  slight_left, slight_right, stop
  Speed: integer 0-255. Camera: 640x480.
Limitations: no arms, no stairs, turns approximate.
Sensors now:
  front: 18.4 cm  WARNING: obstacle ahead
  left:  62.0 cm
  right: 25.3 cm
Scene: "A desk with laptop and water bottle"
Last actions: [forward, forward, slight_left]
State: SEARCHING
---
User said: "find my water bottle"
Respond ONLY with valid JSON.
\end{lstlisting}

\subsection{Task Agent and Episodic Visual Search}
\label{sec:search}

Intent labels are \texttt{COMMAND}, \texttt{QUESTION}, \texttt{TASK}, or \texttt{UNKNOWN}.
\texttt{TASK} plans are JSON step lists (move, speak, \texttt{capture\_check}, \texttt{loop\_while}, wait) capped at 50 steps.
Search uses Algorithm~\ref{alg:search}: capture the frame, get a VLM description, filter by similarity (Eq.~\ref{eq:similarity}), check if target is visible, then select a move based on ultrasonic geometry. This is related to ReAct~\cite{yao2023react} but without an explicit thought trace from the robot.

\begin{algorithm}[t]
\caption{Episodic visual search (VLM-driven)}
\label{alg:search}
\begin{algorithmic}[1]
\REQUIRE target $T$, max steps $N{=}20$
\STATE $\mathcal{M} \leftarrow \emptyset$; $\tau \leftarrow 0.7$; $\mathit{skip} \leftarrow 0$
\FOR{$i = 1$ \TO $N$}
  \STATE $f \leftarrow \text{capture()}$;
         $s \leftarrow \text{VLM.describe}(f)$
  \IF{$\exists\, m \in \mathcal{M}: \text{sim}(s, m) > \tau$}
    \STATE $\mathit{skip} \leftarrow \mathit{skip} + 1$
    \IF{$\mathit{skip} > 5$}
      \STATE $\tau \leftarrow 0.5$
    \ENDIF
    \STATE \textbf{continue}
  \ENDIF
  \STATE $\mathcal{M} \leftarrow \mathcal{M} \cup \{s\}$; $\mathit{skip} \leftarrow 0$
  \IF{$\text{VLM.check}(f, T) = \textsc{yes}$}
    \RETURN \textsc{found}, $s$
  \ENDIF
  \STATE execute(heuristic\_move(sensors)); sleep
\ENDFOR
\RETURN \textsc{not\_found}
\end{algorithmic}
\end{algorithm}

\subsection{Safety Validator (Critic)}
\label{sec:validator}

The critic enforces invariants I$_1$--I$_4$:

\begin{align}
\text{I}_1&: \; r \text{ is a dict with keys \texttt{action}, \texttt{speed}} \label{inv:dict}\\
\text{I}_2&: \; r[\texttt{action}] \in \mathcal{A} \;\text{ else } \safe{} \label{inv:action}\\
\text{I}_3&: \; r[\texttt{speed}] \in \mathbb{Z} \;\text{ after parse else } \safe{} \label{inv:int}\\
\text{I}_4&: \; 0 \leq r[\texttt{speed}] \leq 255 \;\text{ via clamp} \label{inv:range}
\end{align}

\begin{lstlisting}[language=Python, caption={Critic $V$: total, safe output (matches \texttt{safety\_validator.py})},
  label={lst:validator}, float=tp]
SAFE = {"action": "stop", "speed": 0}
def validate(cmd):
    if not isinstance(cmd, dict):  return SAFE
    action = cmd.get("action")
    speed  = cmd.get("speed", 100)
    if action not in VALID_ACTIONS: return SAFE
    try:    speed = int(speed)
    except: return SAFE
    speed = max(0, min(255, speed))
    return {"action": action, "speed": speed}
\end{lstlisting}

\subsection{Firmware Reflex and Pi Cancel}
\label{sec:reflex}

On each ESP32 loop, if \texttt{movingForward} and front distance $d_F < 15$\,cm, motors stop, state becomes \texttt{obstacle}, and telemetry reports \texttt{OBSTACLE}.
Mean measured stop latency is $47 \pm 21$\,ms (Table~\ref{tab:timing}), dominated by the 100\,ms loop period and ultrasonic sampling, not cloud RTT.
The Pi heartbeat (100\,ms period) updates body context and sets a global cancel flag so deliberative states exit cleanly.

\subsection{Seven-State Machine}
\label{sec:fsm}

States: \texttt{IDLE}, \texttt{LISTENING}, \texttt{THINKING}, \texttt{ACTING}, \texttt{SEARCHING}, \texttt{SPEAKING}, \texttt{REPORTING} (Fig.~\ref{fig:statemachine}).
Four Pi threads run concurrently: wake word, ESP32 heartbeat, Telegram poll, and TTS queue.
No cloud call or search step may block reflex handling; cancel is event-driven from telemetry.

\begin{figure}[t]
\centering
\slotfig{0.9\columnwidth}{5.0cm}{statemachine}
\caption{Seven-state machine with global cancel on obstacle detection or stop command.}
\label{fig:statemachine}
\end{figure}

\subsection{Telegram and Presentation Mode}

Telegram accepts text, voice notes (Groq audio transcription), and images (NVIDIA NIM visual analysis) for search-by-photo.
Presentation mode cycles scripted demos with OLED face animations for operator feedback.

% ====================================================================
\section{Evaluation}
\label{sec:results}
% ====================================================================

\subsection{Experimental Protocol}
\label{sec:protocol}

\textbf{Environment.} Single $3\times 4$\,m indoor lab with desks, chairs, and bookshelf; ambient noise $<$40\,dB; local Wi-Fi.
\textbf{Software.} Vision module commit tested May 2025; Groq Llama 3.3 and NVIDIA NIM Llama 3.2 Vision; firmware \texttt{firmware\_phase4.ino}; Pi OS with Python~3.10+.
\textbf{Operators.} Two authors; trials recorded manually in a spreadsheet.
\textbf{Metrics.} \emph{Intent correct}: label matches scripted category; \emph{Action correct}: executed motion or search outcome matches request; \emph{Failure}: unsafe motion, parse failure, or search timeout; \emph{Near-miss}: obstacle within 15\,cm after a validated forward command while still moving.
\textbf{Voice set.} Six paraphrase classes, 10 trials each ($n{=}60$): forward, left turn, stop, scene question, bottle search, non-primitive ``dance.''
\textbf{Search.} Ten trials, five object classes, random placement; success = target reported found within 20 steps.
\textbf{Ablation.} Five scripted commands (forward, turn, stop, describe, search) repeated 10 times per condition; one layer disabled per column in Table~\ref{tab:ablation}.
\textbf{Baselines.} Full stack vs.\ (a)~prompt-only (validator always passes) simulated offline on logged model outputs; (b)~ESP32 \texttt{AUTO\_MODE} fallback when Pi UART is silent (Phase~2 rule-based avoidance, no voice).
\textbf{Ethics.} Voice and images processed via cloud APIs; demonstrations use consenting operators only.

Hardware cost appears in Table~\ref{tab:bom}.

\subsection{Voice Commands}

Table~\ref{tab:voice} summarizes voice results.

\begin{table}[t]
\caption{Voice command evaluation ($n{=}10$ per utterance class).}
\label{tab:voice}
\centering
\renewcommand{\arraystretch}{1.05}
\begin{tabularx}{\columnwidth}{@{}YCC@{}}
\toprule
\textbf{Utterance class} & \textbf{Intent correct} & \textbf{Outcome correct} \\
\midrule
``Go forward'' & 10/10 & 10/10 \\
``Turn left'' & 10/10 & 10/10 \\
``Stop'' & 10/10 & 10/10 \\
``What do you see?'' & 9/10 & 9/10 \\
``Find my bottle'' & 10/10 & 8/10 \\
``Dance for me'' & 8/10 & 7/10 \\
\bottomrule
\end{tabularx}
\end{table}

Macro-averaged precision 0.95, recall 0.93, F1$=$0.94.
Wake phrase \emph{Hey Kanda} (openWakeWord custom model \texttt{hey\_kanda.onnx}): $>$95\% true-positive rate at 1\,m, $<$2\% false-positive rate over 30\,min idle.

\subsection{Visual Object Search}

Eight of ten trials found the target in mean 7.3 steps.
Failures were full occlusion behind furniture.
Episodic memory avoided revisits; $\tau$ adaptation fired three times.

\subsection{Fault Injection}

Table~\ref{tab:faults} lists eight injected faults; all ended in \safe{} or idle within 200\,ms with no sustained forward motion into obstacles.

\begin{table}[t]
\caption{Fault injection ($n{=}1$ deliberate trial per scenario).}
\label{tab:faults}
\centering
\renewcommand{\arraystretch}{1.05}
\begin{tabularx}{\columnwidth}{@{}YY@{}}
\toprule
\textbf{Scenario} & \textbf{Observed behaviour} \\
\midrule
Obstacle in path & ESP32 reflex; stop in 47\,ms (mean, Table~\ref{tab:timing}) \\
Cloud API timeout (15\,s) & Spoken error; return to idle \\
API rate limit & Backoff 5\,s; retry succeeded \\
Camera blocked & Frame rejected ($<$5\,KB); step skipped \\
API key revoked & Validator emits \safe{} each cycle \\
USB unplugged & Pi sends \safe{}; ESP32 telemetry continues \\
Heavy noise & STT failure message; idle \\
Image without target & Search reports not found \\
\bottomrule
\end{tabularx}
\end{table}

\subsection{Ablation}

Table~\ref{tab:ablation} reports failure rates (95\% binomial CI, $n{=}10$ per condition).

\begin{table}[t]
\caption{Ablation: distinct failure modes per removed layer ($n{=}10$, 95\% CI).}
\label{tab:ablation}
\centering
\renewcommand{\arraystretch}{1.05}
\begin{tabularx}{\columnwidth}{@{}YCC@{}}
\toprule
\textbf{Removed layer} & \textbf{Failure rate} & \textbf{Dominant mode} \\
\midrule
Camera scene in prompt & 15\% $\pm$ 7\% & Plan toward unseen obstacle \\
Episodic memory & 60\% $\pm$ 10\% & Revisit prior locations \\
Obstacle warnings in prompt & 20\% $\pm$ 8\% & Forward near wall \\
Action vocabulary in prompt & 30\% $\pm$ 9\% & Non-JSON / bad action \\
ESP32 firmware reflex & 30\% $\pm$ 9\% & Near-miss (see \S\ref{sec:protocol}) \\
\bottomrule
\end{tabularx}
\end{table}

Prompt-only baseline on cached outputs yielded 42\% structurally unsafe commands vs.\ 0\% after $V$ (offline replay of 50 logged responses).
\texttt{AUTO\_MODE} avoided collisions in 10/10 UART-loss trials but could not follow voice goals.

\subsection{Comparison with Prior Systems}

Table~\ref{tab:comparison} contrasts platforms on cost, training, grounding, safety, and modalities.
Among the listed systems, KANDA is the only one we implement with explicit critic--reflex layers on sub-\$100 hardware; larger systems may include implicit filters not shown in publications.

\begin{table}[!t]
\caption{Comparison with embodied LLM robotics systems.}
\label{tab:comparison}
\centering
\scriptsize
\setlength{\tabcolsep}{2pt}
\renewcommand{\arraystretch}{0.95}
\begin{tabularx}{\columnwidth}{@{}YYYYYY@{}}
\toprule
\textbf{Sys.} & \textbf{Cost} & \textbf{Train.} & \textbf{Ground.} & \textbf{Safety} & \textbf{Modal.} \\
\midrule
SayCan & $>$100k & 68k demos & Learned afford. & Afford. filter & Lang. \\
RT-1 & $>$100k & 130k ep. & E2E policy & --- & Vis.+lang. \\
RT-2 & $>$100k & RT-1+web & Fine-tuned VLM & --- & Vis.+lang. \\
PaLM-E & Res. manip. & Multi-task & Sensor tokens & --- & Vis.+lang. \\
Wake et al. & $>$10k & 0-shot & API prompt & --- & Lang. \\
\textbf{KANDA} & \textbf{$<$\$77} & \textbf{0-shot} & \textbf{Body-context} & \textbf{$V$+reflex} & \textbf{V+T+txt} \\
\bottomrule
\end{tabularx}
\end{table}

\subsection{Timing and Cost}

Table~\ref{tab:timing} lists pipeline latencies ($n{=}20$ cycles).
Reflex latency is $\sim$40$\times$ faster than mean intent classification, motivating separate Pi threads and ESP32 reflex logic.

\begin{table}[t]
\caption{End-to-end latencies ($n{=}20$ cycles, local Wi-Fi).}
\label{tab:timing}
\centering
\renewcommand{\arraystretch}{1.05}
\begin{tabularx}{\columnwidth}{@{}YCC@{}}
\toprule
\textbf{Stage} & \textbf{Mean $\pm$ SD} & \textbf{Range} \\
\midrule
Wake word detection & $62 \pm 18$\,ms & 38--95\,ms \\
Intent classification & $1.72 \pm 0.41$\,s & 1.0--2.5\,s \\
Plan generation (6 steps) & $2.89 \pm 0.53$\,s & 2.0--4.0\,s \\
VLM scene description & $2.31 \pm 0.58$\,s & 1.5--3.5\,s \\
Camera capture + encode & $142 \pm 26$\,ms & 98--195\,ms \\
ESP32 JSON parse & $2.1 \pm 0.8$\,ms & 1.2--4.3\,ms \\
Firmware reflex stop & $47 \pm 21$\,ms & 18--92\,ms \\
\bottomrule
\end{tabularx}
\end{table}

\begin{table}[t]
\caption{Bill of materials (2025 retail, USD).}
\label{tab:bom}
\centering
\renewcommand{\arraystretch}{1.05}
\begin{tabularx}{\columnwidth}{@{}YR@{}}
\toprule
\textbf{Component} & \textbf{Cost} \\
\midrule
Raspberry Pi 4 (4\,GB) & \$35 \\
ESP32 DevKit             & \$4 \\
Pi Camera v2.1           & \$10 \\
Motor driver + sensors   & \$5 \\
Display, mic, speaker    & \$10 \\
Chassis, motors, battery & \$13 \\
\midrule
\textbf{Total}           & \textbf{$<$ \$77} \\
\bottomrule
\end{tabularx}
\end{table}

% ====================================================================
\section{Discussion}
\label{sec:discussion}
% ====================================================================

A command safe at validation time may be unsafe after the environment changes; the reflex layer closes that gap.
The critic cannot fix semantic mistakes that remain syntactically valid (e.g., forward toward a distant wall); prompt warnings and scene context reduce but do not eliminate those cases.
Deliberation and reflex operate on different timescales; removing either reintroduces failures shown in Table~\ref{tab:ablation}.

\textbf{Limitations.} Three ultrasonics leave blind spots; no manipulation; no SLAM or metric map; single-room evaluation; cloud dependency and API cost; wake phrase limited to \emph{Hey Kanda} in tests; small $n$ for ablation; no inter-rater labeling.
\textbf{Future work.} On-device Gemma-class model, depth camera, wheel odometry, larger multi-room study, and hardware-in-the-loop logged benchmarks.

% ====================================================================
\section{Conclusion}
\label{sec:conclusion}
% ====================================================================

KANDA demonstrates a low-cost multimodal embodied agent with explicit deliberative--critic--reflex safety.
On our test suite, intent F1 is 0.94, visual search succeeds in 80\% of trials, and injected faults terminate safely.
Ablation confirms that prompt grounding, episodic search memory, syntactic validation, and firmware reflex address largely non-overlapping failure modes.

% ====================================================================
\section*{Acknowledgment}
% ====================================================================

The authors thank the Department of Information Science and Engineering, RV College of Engineering, Bengaluru, for laboratory support and project guidance.

% ====================================================================
\begin{thebibliography}{00}

\bibitem{ahn2022saycan}
M. Ahn et al.,
``Do As I Can, Not As I Say: Grounding Language in Robotic Affordances,''
in \textit{Proc. Conf. Robot Learning (CoRL)}, 2022.

\bibitem{driess2023palme}
D. Driess et al.,
``PaLM-E: An Embodied Multimodal Language Model,''
in \textit{Proc. 40th Int. Conf. Machine Learning (ICML)}, 2023.

\bibitem{brohan2023rt2}
A. Brohan et al.,
``RT-2: Vision-Language-Action Models Transfer Web Knowledge to
Robotic Control,'' \textit{arXiv:2307.15818}, Jul.\ 2023.

\bibitem{brohan2022rt1}
A. Brohan et al.,
``RT-1: Robotics Transformer for Real-World Control at Scale,''
in \textit{Proc. RSS}, 2023. [arXiv:2212.06817]

\bibitem{vemprala2024chatgpt}
S. Vemprala, R. Bonatti, A. Bucker, and A. Kapoor,
``ChatGPT for Robotics: Design Principles and Model Abilities,''
\textit{IEEE Access}, vol.~12, pp.~1--20, 2024.

\bibitem{liang2023code}
J. Liang et al.,
``Code as Policies: Language Model Programs for Embodied Control,''
in \textit{Proc. IEEE ICRA}, 2023.

\bibitem{singh2023progprompt}
I. Singh et al.,
``ProgPrompt: Program Generation for Situated Robot Task Planning,''
in \textit{Proc. IEEE ICRA}, 2023.

\bibitem{huang2023grounded}
W. Huang et al.,
``Grounded Decoding: Guiding Text Generation with Grounded Models for
Robot Control,'' in \textit{Proc. NeurIPS}, 2023.

\bibitem{huang2022zeroshot}
W. Huang, P. Abbeel, D. Pathak, and I. Mordatch,
``Language Models as Zero-Shot Planners,''
in \textit{Proc. ICML}, 2022. [arXiv:2201.07207]

\bibitem{huang2022innermonologue}
W. Huang et al.,
``Inner Monologue: Embodied Reasoning Through Planning with LLMs,''
in \textit{Proc. CoRL}, 2022. [arXiv:2207.05608]

\bibitem{yao2023react}
S. Yao et al.,
``ReAct: Synergizing Reasoning and Acting in Language Models,''
in \textit{Proc. ICLR}, 2023. [arXiv:2210.03629]

\bibitem{wu2023tidybot}
J. Wu et al.,
``TidyBot: Personalized Robot Assistance with Large Language Models,''
in \textit{Proc. IEEE/RSJ IROS}, 2023.

\bibitem{wake2023gptrobot}
N. Wake, A. Kanehira, K. Sasabuchi, J. Takamatsu, and K. Ikeuchi,
``ChatGPT Empowered Long-Step Robot Control in Various Environments,''
in \textit{Proc. IEEE/RSJ IROS}, 2023.

\bibitem{zeng2023socratic}
A. Zeng et al.,
``Socratic Models: Composing Zero-Shot Multimodal Reasoning with Language,''
in \textit{Proc. ICLR}, 2023.

\bibitem{wei2022cot}
J. Wei et al.,
``Chain-of-Thought Prompting Elicits Reasoning in Large Language Models,''
in \textit{Proc. NeurIPS}, vol.~35, 2022.

\bibitem{amodei2016safety}
D. Amodei et al.,
``Concrete Problems in AI Safety,'' \textit{arXiv:1606.06565}, 2016.

\bibitem{brooks1986robust}
R. A. Brooks,
``A Robust Layered Control System for a Mobile Robot,''
\textit{IEEE J. Robotics Autom.}, vol.~2, no.~1, pp.~14--23, 1986.

\bibitem{fox1997dwa}
D. Fox, W. Burgard, and S. Thrun,
``The Dynamic Window Approach to Collision Avoidance,''
\textit{IEEE Robotics Autom. Mag.}, vol.~4, no.~1, pp.~23--33, 1997.

\bibitem{bonomi2012fog}
F. Bonomi, R. Milito, J. Zhu, and S. Addepalli,
``Fog Computing and its Role in the Internet of Things,''
in \textit{Proc. 1st ACM Workshop Mobile Cloud Computing}, pp.~13--16, 2012.

\bibitem{shi2016edge}
W. Shi, J. Cao, Q. Zhang, Y. Li, and L. Xu,
``Edge Computing: Vision and Challenges,''
\textit{IEEE Internet Things J.}, vol.~3, no.~5, pp.~637--646, 2016.

\bibitem{ratcliff1988}
J. W. Ratcliff and D. Metzener,
``Pattern Matching: The Gestalt Approach,''
in \textit{Proc. Dr. Dobb's Journal}, Jul.\ 1988.

\bibitem{openwakeword2024}
D. Kamminga,
``openWakeWord: Open-Source Audio Wake Word Detection,''
GitHub repository, 2024. [Online]. Available: https://github.com/dscripka/openWakeWord

\end{thebibliography}

\end{document}
